% ch07_u01 —— 原书 p.236--p.245（7.1.1 -- 7.3.3，第 7 章开头）
% 本单元自第 7 章标题起，至 7.3.3 末句止；尾部接 p.246。

\bdchap{ch:appstack}{附录：代数栈上的 D-模理论与 Hecke 模式}

\SourcePage{236}

\bdsec{7.1. 引言}

\bdp{7.1.1} 本节的中心目标是陈述一个一般的 Hecke 模式，它用于主定理的证明。
其（非扭）有限维版本如下。设 $G$ 为代数群，$K\subset G$ 为代数子群，
$\mathfrak{g}$ 为 $G$ 的李代数，$Y$ 为带 $G$-作用的光滑簇。记
$\mathcal{H}:=D(K\backslash G/K)$ 为栈 $K\backslash G/K$ 的 $\mathcal{D}$-模导出范畴。
类似地有导出范畴 $D(K\backslash Y)$，以及 $(\mathfrak{g},K)$-模的范畴
$\mathcal{M}(\mathfrak{g},K)$ 的导出范畴 $D(\mathfrak{g},K)$。于是有如下
「Hecke 模式」：
\begin{enumerate}[label=(\alph*),leftmargin=2.2em]
\item $\mathcal{H}$ 是幺半三角范畴；
\item $D(K\backslash Y)$ 是 $\mathcal{H}$-模；
\item $D(\mathfrak{g},K)$ 是 $\mathcal{H}$-模；
\item 标准函子
\[
{}^{L}\Delta:D(\mathfrak{g},K)\longrightarrow D(K\backslash Y),
\qquad
R\Gamma:D(K\backslash Y)\longrightarrow D(\mathfrak{g},K)
\]
是 $\mathcal{H}$-模的态射。
\end{enumerate}

这里 ${}^{L}\Delta$ 与 $R\Gamma$ 是 1.2.4 中函子 $\Delta$ 与 $\Gamma$ 的导出版本。
$\mathcal{H}$ 上的张量积，以及 (b)、(c) 中的 $\mathcal{H}$-作用，
都是适当的「卷积」函子 $\circledast$。例如考察 $K=\{1\}$ 的情形。
记 $\delta_g$ 为 $g\in G$ 处的 $\delta$-函数对应的 $\mathcal{D}$-模，则有
$\delta_{g_1}\circledast\delta_{g_2}=\delta_{g_1g_2}$。对 $Y$ 上的
$\mathcal{D}$-模 $M$，$\delta_g\circledast M$ 是 $M$ 的 $g$-平移；对
$\mathfrak{g}$-模 $V$，$\delta_g\circledast V$ 就是 $V$ 配上经
$\operatorname{Ad}_g$ 扭过的 $\mathfrak{g}$-作用。$M$ 上的 $\mathcal{D}$-模结构
把无限接近的 $g$ 所对应的 $\delta_g\circledast M$ 典范地等同起来；类似地，
$V$ 上的 $\mathfrak{g}$-作用把这些 $\delta_g\circledast V$ 等同起来。
这就允许对 $G$ 上任意的 $\mathcal{D}$-模定义卷积函子。

\bdp{7.1.2} Hecke 函子的精确构造需要某种栈上的 $\mathcal{D}$-模形式化。例如，
人们需要光滑栈 $\mathcal{Y}$ 的 $\mathcal{D}$-模%
\SourcePage{237}%
导出范畴 $D(\mathcal{Y})$ 的定义（它可能与 $\mathcal{Y}$ 上 $\mathcal{D}$-模范畴的
导出范畴并不相同！）。似乎没有现成的参考文献。例外是 $\mathcal{Y}$ 为轨道栈的
情形，即光滑簇对仿射群作用的商；[BL]、[Gi87] 处理过它，但方式不太便于
Hecke 函子的应用。因此我们只好补上一些一般废话，以免沉没。

我们仿照 Kapranov [Kap91] 与 Saito [Sa89]，在 7.2 中从如下典范等价开始：
$\mathcal{D}$-模的导出范畴与 $\Omega$-模的导出范畴之间的等价
（这里 $\Omega$ 是微分形式的 DG 代数），它把 $\mathcal{D}$-模等同于其 de Rham 复形。
处理栈时，$\Omega$-模更易驾驭：原因在于 $\Omega$ 是光滑拓扑上的环层，
而 $\mathcal{D}$ 不是。在重要的特殊情形——栈的对角态射为仿射——这一「超」%
\footnote{这是某位数学家对 Mary Poppins 生造的词
「supercalifragelistic-expialidocious」的缩写。
% [Tutor Note: 原书 p.237 脚注把该生造词印作「supercalifragelistic-」，通行拼法为
% 「supercalifragilisticexpialidocious」；底本如此，译文照录。]
}%
格式尤为方便。此时可以（见 7.3）直接用「整体」$\Omega$-复形来定义 $\mathcal{D}$-模
导出范畴。在 7.5 中，经过 7.4 的一般同调代数离题之后，我们给出 $\mathcal{D}$-模
导出范畴的一个「局部」定义，它对任意的光滑栈都适用。

7.6 解释 Hecke 模式的 (a)、(b) 两部分；我们还证明，轨道栈的
$\mathcal{D}$-模导出范畴等价于 [BL]、[Gi87] 中的等变导出范畴。7.7 描述
Harish-Chandra 模的类似「超」格式；作为额外收获，我们在 7.7.12 得到 [BL]
主要结果的一个简单证明。Hecke 模式的 Harish-Chandra 部分 (c)、(d) 在 7.8 中处理。
正文所需的带额外对称与参数的版本在 7.9 中给出。在转入无穷维设定之前，
我们在 7.10 讨论 $\mathcal{D}$-模的晶体方法，处理奇异空间时它尤为方便
（本节要归功于 1980 年与 J. Bernstein 的讨论）。7.11 与 7.12 两节包含关于
归纳概形、Mittag-Leffler 模的一些基本材料，%
\SourcePage{238}%
以及形式光滑归纳概形上的 $\mathcal{D}$-模。7.13 节综述 BRST 约化。Hecke 模式
(c)、(d) 两部分的无穷维版本在 7.14 中。最后，7.15 证明仿射旗簇上的正扭
$\mathcal{D}$-模本质上等同于小于临界水平的仿射 Kac–Moody 李代数的表示。
在 $\mathcal{D}$-模沿 Schubert 分层光滑的特殊情形，Kashiwara 与 Tanisaki [KT95]
得到了类似的结果（[KT95] 的作者不使用归纳概形上 $\mathcal{D}$-模的语言）。
我们还把相应的 de Rham 上同调群与 BRST 上同调群等同起来。

我们对 $\mathcal{D}$-模理论的叙述相当不完整；基本上只处理本文正文用到的题材。
例外是 7.4、7.5 两节（栈 $\operatorname{Bun}_G$ 落入 7.3 的形式化）、7.10 节
（我们遇到的奇异空间都是仿射 Grassmann 流形上的分层，故可用 7.11）以及 7.15 节
（它收进来只为博读者一乐）。

回忆：$\mathcal{M}^{\ell}(X)$（相应地 $\mathcal{M}^{r}(X)$）表示光滑簇 $X$ 上
左（相应地右）$\mathcal{D}$-模的范畴；我们常把这些范畴等同起来，并记作
$\mathcal{M}(X)$。若 $F$ 是复形，则记 $F^{\cdot}$ 为相应的分次对象
（忘掉微分）。

\bdsec{7.2. $\mathcal{D}$- 与 $\Omega$-模}

\bdp{7.2.1} 设 $X$ 为光滑代数簇%
\footnote{或者更一般地，$\mathbb{C}$ 上光滑拟紧的代数空间，其对角态射
$X\to X\times X$ 为仿射。本节（7.2.10 除外）的构造与陈述都是局部的，
故对任意光滑代数空间都有意义。对 $X$ 的这一要求是为了保证我们定义的导出范畴
满足适当的局部-整体（下降）性质。我们在 7.5 中于栈的更一般设定下讨论此事。}%
。记 $\Omega_X$ 为 $X$ 上微分形式的 DG 代数。则 $(X,\Omega_X)$ 是 DG 环化空间，
于是有 $\Omega_X$-复形的范畴（$:=$ DG $\Omega_X$-模）。$\Omega_X$-复形
$F=(F^{\cdot},d)$ 称为拟凝聚的，如果各 $F^i$ 都是拟凝聚的 $\mathcal{O}_X$-模；
拟凝聚的 $\Omega_X$-复形通常称为 $X$ 上的 \emph{$\Omega$-复形}。记%
\SourcePage{239}%
为 $X$ 上 $\Omega$-复形的 DG 范畴 $C(X,\Omega)$。这是张量 DG 范畴。

\emph{注.} 对 $\Omega_X$-复形 $F$，微分 $d:F^{\cdot}\to F^{\cdot+1}$
是阶 $\le1$ 的微分算子，其象征等于乘积映射
$\Omega_X^1\otimes F^{\cdot}\to F^{\cdot+1}$。由此可见，$F^{\cdot}$ 上的
$\Omega_X^{\cdot}$-模结构可以由 $\mathcal{O}_X$-模结构与 $d$ 重构出来。事实上，
忘掉 $\Omega_X^{\ge1}$-作用，就把 $C(X,\Omega)$ 与如下复形 $(F^{\cdot},d)$
的范畴等同起来：各 $F^{\cdot}$ 是拟凝聚 $\mathcal{O}_X$-模，$d$ 是阶 $\le1$
的微分算子。

\bdp{7.2.2} 设 $C(X,\mathcal{D}):=C(\mathcal{M}^{r}(X))$ 为 $X$ 上右
$\mathcal{D}$-模的复形（\emph{右 $\mathcal{D}$-复形}，或简称
\emph{$\mathcal{D}$-复形}）的 DG 范畴，$K(X,\mathcal{D})$ 为相应的同伦范畴。
我们有一对伴随 DG 函子
\begin{equation}\tag{306}
\mathcal{D}:C(X,\Omega)\longrightarrow C(X,\mathcal{D}),
\qquad
\Omega:C(X,\mathcal{D})\longrightarrow C(X,\Omega)
\end{equation}
其定义如下。记 $DR_X$ 为 $\mathcal{D}_X$ 视为左 $\mathcal{D}$-模时的 de Rham
复形，故 $DR_X=\Omega_X^{\cdot}\otimes_{\mathcal{O}_X}\mathcal{D}_X$。
这是一个带 $\mathcal{D}_X$ 右作用的 $\Omega$-复形。现在，对 $\Omega$-复形 $F$
与右 $\mathcal{D}$-复形 $M$，有
\begin{equation}\tag{307}
\mathcal{D}F=F\otimes_{\Omega_X}DR_X,
\qquad
\Omega M:=\operatorname{Hom}_{\mathcal{D}_X}(DR_X,M).
\end{equation}
伴随性质是显然的。

\bdp{7.2.3} \emph{注.} (i) 有
$\mathcal{D}F^{\cdot}=F^{\cdot}\otimes_{\mathcal{O}_X}\mathcal{D}_X
=\operatorname{Diff}(\mathcal{O},F^{\cdot})$；微分
$d_{\mathcal{D}F}:\mathcal{D}F^{\cdot}\to\mathcal{D}F^{\cdot+1}$ 把微分算子
$a:\mathcal{O}_X\to F^{\cdot}$ 送到复合 $d\cdot a$。$\Omega$-复形 $\Omega M$
就是 $M$ 的 de Rham 复形，其中
$(\Omega M)^i=\bigoplus_{a-b=i}M^a\otimes\Lambda^b\Theta_X$。

(ii) $X$ 上左 $\mathcal{D}$-模的范畴 $\mathcal{M}^{\ell}(X)$ 按通常方式
（张量积取在 $\mathcal{O}_X$ 上）是张量范畴，故左 $\mathcal{D}$-复形的范畴
$C(\mathcal{M}^{\ell}(X))$ 是张量 DG 范畴。DG 函子 $\Omega:$%
\SourcePage{240}%
$C(\mathcal{M}^{\ell}(X))\to C(X,\Omega)$，它把左 $\mathcal{D}$-复形 $N$ 送到其
de Rham 复形 $(\Omega N)^{\cdot}=\Omega_X^{\cdot}\otimes_{\mathcal{O}_X}N$，
是张量函子。

(iii) DG 范畴 $C(X,\Omega)$ 与 $C(X,\mathcal{D})$ 都是张量 DG 范畴
$C(\mathcal{M}^{\ell}(X))$ 上的模。函子 $\mathcal{D}$ 与 $\Omega$ 都是
$C(\mathcal{M}^{\ell}(X))$-模的态射。

\bdp{7.2.4} \emph{引理.} 对任意 $\mathcal{D}$-复形 $M$，典范态射
$\mathcal{D}\Omega M\to M$ 是拟同构。

\emph{证明.} 令
\[
V_j^{i}:=\bigoplus_{\substack{a-b=i\\ b+c=j}}
M^a\otimes\Lambda^b\Theta_X\otimes\mathcal{D}_X^{\leq c}
\subset(\mathcal{D}\Omega M)^i.
\]
则 $V_{*}$ 是 $\mathcal{D}\Omega M$ 的递增滤链，各项为 $\mathcal{O}$-子复形，
且 $V_0\xrightarrow{\sim}M$，而对 $i\ge1$，$V_i/V_{i-1}$ 无环
（因为 $V_i/V_{i-1}$ 是 $M$ 与 $\Theta_X$ 的第 $i$ 个 Koszul 复形之张量积）。
\hfill$\square$

\bdp{7.2.5} 对 $\Omega$-复形 $F$，置 $H_{\mathcal{D}}^{\cdot}F=H^{\cdot}\mathcal{D}F$。
于是 $H_{\mathcal{D}}$ 是 $K(X,\Omega)$ 上的上同调函子，取值于阿贝尔范畴
$\mathcal{M}^{r}(X)$。$\Omega$-复形态射 $\phi:F_1\to F_2$ 称为
\emph{$\mathcal{D}$-拟同构}，如果 $\mathcal{D}$-复形态射
$\mathcal{D}\phi:\mathcal{D}F_1\to\mathcal{D}F_2$ 是拟同构，亦即
$H_{\mathcal{D}}^{\cdot}F_1\to H_{\mathcal{D}}^{\cdot}F_2$ 是同构。我们有如下
简单性质（用 7.2.4 证明 (ii)、(iii)）：
\begin{enumerate}[label=(\roman*),leftmargin=2.2em]
\item 若 $\phi$ 是 $\mathcal{D}$-拟同构，$N$ 是作为 $\mathcal{O}$-模平坦的左
$\mathcal{D}$-模，则 $\phi\otimes\operatorname{id}_N:F_1\otimes N\to F_2\otimes N$
是 $\mathcal{D}$-拟同构。
\item 典范态射 $\alpha_F:F\to\Omega\mathcal{D}F$ 是 $\mathcal{D}$-拟同构。
\item $\Omega$ 把拟同构送到 $\mathcal{D}$-拟同构。
\end{enumerate}

下面的引理在后文中不会用到；读者可以跳过。若 $\Omega$-复形态射
$\phi:F_1\to F_2$ 是层复形（取值于向量空间）的拟同构，则称之为
\emph{朴素的拟同构}。

\SourcePage{241}

\bdp{7.2.6} \emph{引理.} (i) 任何 $\mathcal{D}$-拟同构都是朴素的拟同构。

(ii) 如上态射 $\phi$ 是 $\mathcal{D}$-拟同构，当且仅当对任意下有界的局部自由
$\Omega$-模复形 $A$，态射 $\phi\otimes\operatorname{id}_A:F_1\otimes A\to F_2\otimes A$
都是朴素的拟同构。

(iii) 假设或者 $\Omega^{\ge1}F_i^{\cdot}=0$（即微分是 $\mathcal{O}$-线性的），
或者各 $F_i^{\cdot}$ 有界且 $\mathcal{O}$-凝聚。则任何朴素的拟同构 $\phi$ 都是
$\mathcal{D}$-拟同构。对一般的 $\Omega$-复形，这未必成立。

\emph{证明.} (i) 对任意 $\Omega$-复形 $F$，典范态射
$\alpha_F:F\to\Omega\mathcal{D}F$ 是朴素的拟同构。由于 $\Omega$ 把
$\mathcal{D}$-复形的拟同构送到朴素的拟同构，$\Omega(\mathcal{D}\phi)$ 是朴素的
拟同构。现在由 $\alpha_{F_2}\phi=\Omega((\mathcal{D}\phi)\alpha_{F_1}$ 即得结论。
% [Tutor Note: 原书 p.241 此处 $\Omega((\mathcal{D}\phi)\alpha_{F_1}$ 缺一个右括号，
% 底本如此，译文照录。]

(ii) 为证「若」的陈述，只需取 $A=DR_X$。反之，设 $\phi$ 是
$\mathcal{D}$-拟同构。$A$ 上有下有界的递增滤链 $A_i$，使得 $\cup A_i=A$，
且每个 $\operatorname{gr}_iA$ 是生成元位于 $i$ 次处的局部自由 $\Omega_X$-模
（取 $A_i:=\Omega_X\cdot A_{\le i}$）。于是，若所有
$\phi\otimes\operatorname{id}_{\operatorname{gr}_iA}$ 都是朴素的拟同构，
则 $\phi\otimes\operatorname{id}_A$ 也是。这样便可假设 $A$ 是生成元位于固定次数
（比如 0 次）的局部自由 $\Omega_X$-模，即 $A=\Omega N$，其中 $N$ 是作为
$\mathcal{O}$-模局部自由的左 $\mathcal{D}$-模。此时
$\phi\otimes\operatorname{id}_A=\phi\otimes\operatorname{id}_N$，
由 7.2.5 的 (i) 即得结论。

(iii) $\mathcal{O}$-线性的情形显然（因为此时
$\mathcal{D}F=F\otimes_{\mathcal{O}_X}\mathcal{D}_X$）。$\mathcal{O}$-凝聚的
情形由下面应用于 $\mathcal{D}\phi$ 的子引理推出（注意由 7.2.5 的性质 (ii)，
$\mathcal{D}F$ 在 $x$ 处的纤维与 $R\Gamma_x(X,F)$ 一致）。

\emph{子引理.} 设 $\psi:M_1\to M_2$ 是 $X$ 上凝聚 $\mathcal{D}$-模的有限复形
之间的态射。假设对任意 $x\in X(\mathbb{C})$，相应的纤维%
\footnote{这里当然是在通常的导出范畴意义下考虑 $\mathcal{O}$-模纤维。
% [Tutor Note: 原书 p.241 脚注印作「the $\mathcal{O}$-moduli fibers」，其中
% 「moduli」当为「module」；底本如此，译文按原意译出。]
}%
态射 $M_{1x}\to M_{2x}$ 都是拟同构。则 $\psi$ 是拟同构。

\SourcePage{242}

\emph{子引理的证明.} 置 $C=\operatorname{Cone}(\psi)$；记 $Y$ 为 $H^{\cdot}(C)$
的支撑。假设 $\psi$ 不是拟同构，即 $Y$ 非空。必要时限制 $X$，可设 $Y$ 是 $X$
的光滑子簇，且凝聚 $\mathcal{D}_Y$-模
$P^{\cdot}:=i_Y^{!}H^{\cdot}(C)=H^{\cdot}i_Y^{!}(C)$ 作为 $\mathcal{O}_Y$-模
自由。由于对 $x\in Y$ 有 $H^{\cdot}(C_x)=P_x^{\cdot+n}$，其中 $n$ 是 $Y$ 在
$X$ 中的余维，故 $P^{\cdot}=0$，矛盾。

为得到朴素拟同构而非 $\mathcal{D}$-拟同构的例子，只需找到非零 $\mathcal{D}$-模
$M$，使 $\Omega M$ 是层的无环复形。取 $M$ 为常数层，其值为 $\mathcal{D}_X$-模，
等于微分算子环（在 $X$ 的一般点处）的分式域。\hfill$\square$

\bdp{7.2.7} 由于 $H_{\mathcal{D}}$ 是上同调函子，$\mathcal{D}$-拟同构在
$C(X,\Omega)$ 的同伦范畴中构成局部化族。因此相应的局部化 $D(X,\Omega)$ 是
三角范畴（见 [Ve]）；我们称之为 $\Omega$-复形的 \emph{$\mathcal{D}$-导出范畴}。
函子 $\mathcal{D}$、$\Omega$ 给出三角范畴的互逆等价
\begin{equation}\tag{308}
\mathcal{D}:D(X,\Omega)\longrightarrow D(X,\mathcal{D}),
\qquad
\Omega:D(X,\mathcal{D})\longrightarrow D(X,\Omega).
\end{equation}
这里 $D(X,\mathcal{D})=D\mathcal{M}^{r}(X)$。我们常把这样等同起来的这些三角范畴
记为 $D(X)$。也可以考虑有界导出范畴。

\emph{注.} 对下有界的单射 $\mathcal{D}$-模复形 $M$，相应的 $\Omega$-复形
$\Omega M$ 是单射的。因此同伦范畴 $K^{+}(X,\Omega)$ 有许多单射对象。

\bdp{7.2.8} 设 $f:Y\to Z$ 是光滑簇之间的态射。它给出 DG 环化空间的态射
$f_{\Omega}:(Y,\Omega_Y)\to(Z,\Omega_Z)$。于是有相应的 DG 函子
$f_{\Omega}^{\cdot}:C(Z,\Omega)\to C(Y,\Omega)$，
$f_{\cdot}=f_{\Omega\cdot}:C(Y,\Omega)\to C(Z,\Omega)$。
我们先考察拉回函子。

\SourcePage{243}

对左 $\mathcal{D}$-模有通常的拉回函子
$f^{\dagger}:\mathcal{M}^{\ell}(Z)\to\mathcal{M}^{\ell}(Y)$，
$f^{\dagger}(N)=\mathcal{O}_Y\otimes_{f^{-1}\mathcal{O}_Z}f^{-1}N$。有
$\Omega f^{\dagger}(N)=f_{\Omega}^{\cdot}(\Omega N)$。可以把左 $\mathcal{D}$-模
换成右 $\mathcal{D}$-模%
\footnote{用标准等价
$\mathcal{M}^{\ell}(Z)\xrightarrow{\sim}\mathcal{M}(Z)$，
$N\mapsto N\otimes\omega_Z$。}%
，并考虑相应的函子
$f^{\dagger}:\mathcal{M}^{r}(Z)\to\mathcal{M}^{r}(Y)$；此时
$f_{\Omega}^{\cdot}(\Omega M)=\Omega f^{\dagger}M[-\dim Y/Z]$。

若 $f$ 光滑，则对任意 $F\in C(Z,\Omega)$ 有
$H_{\mathcal{D}}^{\cdot}f_{\Omega}^{\cdot}F
=f^{\dagger}H_{\mathcal{D}}^{-\dim Y/Z}F$。于是 $f_{\Omega}^{\cdot}$ 保持
$\mathcal{D}$-拟同构，从而有函子
$f_{\Omega}^{\cdot}:D(Z,\Omega)\to D(Y,\Omega)$。伴随态射
$\mathcal{D}f_{\Omega}^{\cdot}(\Omega M)\to f^{\dagger}M[-\dim U/X]$ 是拟同构。
% [Tutor Note: 原书 p.243 此处写作 $[-\dim U/X]$，而上一式用的是 $\dim Y/Z$；
% 底本如此，译文照录。]

\bdp{7.2.9} \emph{引理.} $\Omega$-复形是光滑拓扑下的局部对象，即拉回函子使
$C(U,\Omega)$（$U\in X_{sm}$）成为 $X$ 的光滑拓扑上的 DG 范畴层。
$\mathcal{D}$-拟同构的概念在 $X_{sm}$ 上是局部的。\hfill$\square$

\bdp{7.2.10} 让我们回到 7.2.8 的情形，考察 DG 函子
$f_{\cdot}:C(Y,\Omega)\to C(Z,\Omega)$。右导出函子
$Rf_{\cdot}:D(Y,\Omega)\to D(Z,\Omega)$ 的定义是正确的。事实上，设
$U_{\cdot}$ 是 $Y$ 的（有限）仿射覆盖（平展的或 Zariski 的均可）。对
$F\in C(Y,\Omega)$，设 $F\to C(F)$ 为 $F$ 相应的 Čech 消解。则
\footnote{例如可由 7.3.9 之后的注推出。}
$f_{\cdot}C(F)\xrightarrow{\sim}Rf_{\cdot}F$。

我们把相应的函子 $D(Y)\to D(Z)$ 记为 $f_{*}$。它与通常的 $\mathcal{D}$-模前推
函子一致。事实上，对 $Y$ 上的 $\mathcal{D}$-复形 $M$ 有
\[
\mathcal{D}f_{\cdot}\Omega M=f_{\cdot}(\Omega M\otimes f^{\dagger}\mathcal{D}_Z)
=f_{\cdot}(\mathcal{D}(\Omega M)\otimes_{\mathcal{D}_Y}f^{\dagger}\mathcal{D}_Z).
\]
由于 $f^{\dagger}\mathcal{D}_Z$ 是平坦的 $\mathcal{O}_Y$-模，而
$\mathcal{D}(\Omega M)$ 是 $M$ 的消解，可见
$\mathcal{D}(\Omega M)\otimes_{\mathcal{D}_Y}f^{\dagger}\mathcal{D}_Z
=M\otimes^{L}_{\mathcal{D}_Y}f^{\dagger}\mathcal{D}_Z$。于是
$f_{*}M=Rf_{\cdot}(M\otimes^{L}_{\mathcal{D}_Y}f^{\dagger}\mathcal{D}_Z)$，证毕。

我们留给读者验证：$Rf_{\cdot}$ 与 $f$ 的复合相容，即典范态射
$R(fg)_{\cdot}\to Rf_{\cdot}Rg_{\cdot}$ 是同构%
\footnote{见 7.3.10(ii) 中该陈述在更一般情形下的证明。}%
，并且这一等同 $(fg)_{*}=f_{*}g_{*}$ 与 $\mathcal{D}$-模理论中的标准等同一致。

\SourcePage{244}

\bdp{7.2.11} 对 $Y$ 上的 $\mathcal{D}$-复形 $M$，记
$M_{\mathcal{O}}\in D(Y,\mathcal{O})$ 为把同一个 $M$ 视为
$\mathcal{O}^{!}$-模的复形。有一个典范的\emph{积分}态射
\begin{equation}\tag{309}
i_f:Rf_{\cdot}(M_{\mathcal{O}})\longrightarrow(f_{*}M)_{\mathcal{O}}
\end{equation}
它位于 $D(Y,\mathcal{O})$ 中，定义如下。只需定义态射
$i_f:f_{\cdot}(M_{\mathcal{O}})\to\mathcal{D}(f_{\cdot}\Omega M)$。现在 $i_f$
是如下复合：
\[
f_{\cdot}(M_{\mathcal{O}})\longrightarrow[\mathcal{D}(f_{\cdot}(M_{\mathcal{O}}))]_{\mathcal{O}}
\longrightarrow[\mathcal{D}(f_{\cdot}\Omega M)]_{\mathcal{O}},
\]
其中箭头来自典范态射 $N\to(\mathcal{D}N)_{\mathcal{O}}$
（取 $N=f_{\cdot}(M_{\mathcal{O}})$）与 $M_{\mathcal{O}}\to\Omega M$。换言之，
$i_f$ 由对显然的态射
$M_{\mathcal{O}}\to(M\otimes^{L}_{\mathcal{D}_Y}f^{\dagger}\mathcal{D}_Z)_{\mathcal{O}}$
施用 $Rf_{\cdot}$ 而得。
% [Tutor Note: 原书 p.244 此处说 $i_f$ 位于 $D(Y,\mathcal{O})$ 中，但按定义
% $i_f$ 取值于 $D(Z,\mathcal{O})$；底本如此，译文照录。]

我们留给读者验证 $i_f$ 与 $f$ 的复合相容。

\bdsec{7.3. 光滑栈上的 $\mathcal{D}$-模理论 I}

% [Tutor Note: 原书 p.244 的 7.3 节标题是行内（run-in）粗斜体，
% 与正文同一行开始；译文按全书的节标题体例单独成行，编号与标题文字不变。]

我们为满足下文条件（310）的光滑栈建立基本的 $\mathcal{D}$-模形式化。
在 7.3.12 中我们修改定义，使得可以去掉拟紧性假设。任意的光滑栈将在 7.5 中处理。

\bdp{7.3.1} 设 $\mathcal{Y}$ 为光滑拟紧代数栈。假设它满足如下条件%
\footnote{这一条件是为了保证我们定义的范畴 $D(\mathcal{Y})$ 具有右局部-整体性质，
见 7.5.3。7.3.1--7.3.3 的构造对任意光滑代数栈都有意义。}%
：
\begin{equation}\tag{310}
\text{对角态射 }\mathcal{Y}\to\mathcal{Y}\times\mathcal{Y}\text{ 是仿射的。}
\end{equation}
等价地说，这意味着存在光滑仿射满射 $U\to\mathcal{Y}$，使 $U$ 是仿射概形。
换言之，把光滑代数簇 $X$ 模去光滑群胚 $Q$%
\footnote{$Q=X\times_{\mathcal{Y}}X$。}%
的作用，所得商即 $\mathcal{Y}$，且结构态射 $Q\to X\times X$ 是仿射的。

\SourcePage{245}

注意 $\Omega(U)$（$U\in\mathcal{Y}_{sm}$）构成 $\mathcal{Y}_{sm}$ 上的 DG 代数层
$\Omega_{\mathcal{Y}}$。$\mathcal{Y}$ 上的 \emph{$\Omega$-复形}是指这样的
DG $\Omega_{\mathcal{Y}}$-模：它作为 $\mathcal{O}_{\mathcal{Y}}$-模是拟凝聚的。
记 $\mathcal{Y}$ 上 $\Omega$-复形的 DG 范畴为 $C(\mathcal{Y},\Omega)$。

\emph{注.} 范畴 $C(U,\Omega)$（$U\in\mathcal{Y}_{sm}$）构成 $\mathcal{Y}_{sm}$
上的 DG 范畴层 $C(\mathcal{Y}_{sm},\Omega)$（见 7.2.9），而 $\mathcal{Y}$ 上的
$\Omega$-复形就是 $C(\mathcal{Y}_{sm},\Omega)$ 的笛卡尔截面。等价地，
$\mathcal{Y}$ 上的 $\Omega$-复形就是 $X$ 上 $Q$-等变的 $\Omega$-复形。

\bdp{7.3.2} 回忆：$\mathcal{D}$-模的范畴 $\mathcal{M}(U)$（$U\in\mathcal{Y}_{sm}$）
构成 $\mathcal{Y}_{sm}$ 上的阿贝尔范畴层，而 $\mathcal{Y}$ 上的
$\mathcal{D}$-模范畴 $\mathcal{M}(\mathcal{Y})$ 即其笛卡尔截面范畴。由 7.2.8 有
典范的上同调函子
$H_{\mathcal{D}}^{\cdot}:C(\mathcal{Y},\Omega)\to\mathcal{M}(\mathcal{Y})$，
$H_{\mathcal{D}}^{\cdot}(F)_U:=H_{\mathcal{D}}^{\cdot+\dim U/\mathcal{Y}}(F_U)$。
$\Omega$-复形的态射若诱导各 $H_{\mathcal{D}}^{\cdot}$ 的同构，则称为
$\mathcal{D}$-拟同构。把 $\Omega$-复形的同伦范畴对 $\mathcal{D}$-拟同构作局部化，
得到三角范畴 $\mathcal{D}(\mathcal{Y})=\mathcal{D}(\mathcal{Y},\Omega)$。
相应地也有有界导出范畴。

有一个全忠实嵌入 $\mathcal{M}(\mathcal{Y})\hookrightarrow\mathcal{D}(\mathcal{Y})$，
它把 $\mathcal{Y}$ 上的 $\mathcal{D}$-模 $M$ 送到其 de Rham 复形 $\Omega M$，
$(\Omega M)_U:=\Omega M_U[-\dim U/\mathcal{Y}]$。有
$H_{\mathcal{D}}^0\Omega M=M$，且对 $a\neq0$ 有 $H_{\mathcal{D}}^a\Omega M=0$。
容易看出，$\Omega$ 把 $\mathcal{M}(\mathcal{Y})$ 与 $\mathcal{D}(\mathcal{Y})$ 中
由那些满足「对 $a\neq0$ 有 $H_{\mathcal{D}}^a(F)=0$」的 $\Omega$-复形 $F$ 组成的
全子范畴等同起来。

\bdp{7.3.3} \emph{例.} 记 $\Omega\mathcal{D}_{\mathcal{Y}}$ 为 $\mathcal{Y}$ 上
由 $\Omega\mathcal{D}_{\mathcal{Y}U}:=\Omega_{U/\mathcal{Y}}[\dim\mathcal{Y}]$
定义的 $\Omega$-复形。注意对 $a>0$ 有
$H_{\mathcal{D}}^a(\Omega\mathcal{D}_{\mathcal{Y}})=0$。若 $\mathcal{Y}$ 是良栈，
则我们的 $\Omega$-复形属于 $\mathcal{M}(\mathcal{Y})$ 的本质像；相应的
$\mathcal{D}$-模 $\mathcal{D}_{\mathcal{Y}}
=H_{\mathcal{D}}^0(\Omega\mathcal{D}_{\mathcal{Y}})$ 与 1.1.3 中的左
$\mathcal{D}$-模 $\mathcal{D}_{\mathcal{Y}}$ 一致。更一般地，对 $\mathcal{Y}$ 上
任意 $\mathcal{O}$-模 $P$，有 $\Omega$-复形
$\Omega(\mathcal{D}_{\mathcal{Y}}\otimes P)$，其中
\[
\Omega(\mathcal{D}_{\mathcal{Y}}\otimes P)_U:=
\Omega_{U/\mathcal{Y}}\otimes_{\mathcal{O}_{\mathcal{Y}}}P_U[\dim\mathcal{Y}].
\]
若 $\mathcal{Y}$ 是良栈且 $P$ 局部自由，则我们的 $\Omega$-复形落在
$\mathcal{M}(\mathcal{Y})$ 中，且等于左 $\mathcal{D}$-模
$\mathcal{D}_{\mathcal{Y}}\otimes P
=\mathcal{D}_{\mathcal{Y}}\otimes_{\mathcal{O}_{\mathcal{Y}}}P$。

% 尾部接 p.246（本单元末句完结于 D_𝒴⊗P = D_𝒴⊗_{O_𝒴}P，7.3.3 至此结束）。
